\documentclass[11pt,a4paper]{article}
\usepackage[utf8]{inputenc}
\usepackage[T1]{fontenc}
\usepackage[spanish,es-minimal,es-noshorthands]{babel}
\AtBeginDocument{\renewcommand{\sin}{\operatorname{sen}}\renewcommand{\arcsin}{\operatorname{arcsen}}}
\usepackage{lmodern}
\usepackage[sfdefault,lining]{FiraSans}
\usepackage[a4paper,top=2.3cm,bottom=2.2cm,left=1.8cm,right=1.8cm,headheight=1.3cm,headsep=0.4cm,footskip=1cm]{geometry}
\usepackage{amsmath,amssymb}
\usepackage{icomma}
\usepackage{graphicx}
\usepackage[table]{xcolor}
\usepackage{tikz}
\usepackage{fontawesome5}
\usepackage{ccicons}
\usepackage{fancyhdr}
\usepackage{lastpage}
\usepackage{tasks}
\usepackage{needspace}
\usepackage[most]{tcolorbox}
\usepackage{array}
\usepackage[hidelinks]{hyperref}
\AtBeginDocument{\hypersetup{pdftitle={Ficha de ejercicios de la unidad 7: Geometría básica del plano. 1º ESO},pdfsubject={Matemáticas 1º ESO},pdfkeywords={punto, recta, semirrecta, segmento, rectas perpendiculares, mediatriz, ángulos, bisectriz, sistema sexagesimal, complementarios, suplementarios, polígonos, diagonales, triángulos, cuadriláteros, polígonos regulares}}}

\definecolor{navy}{HTML}{1B2748}
\definecolor{accent}{HTML}{E8603A}
\definecolor{teal}{HTML}{1F9E90}
\definecolor{soft}{HTML}{EEF1F8}
\definecolor{muted}{HTML}{5D6685}

\setlength{\parindent}{0pt}
\setlength{\parskip}{3pt}
\settasks{label=\alph*), label-format=\bfseries\color{navy}, label-width=1.3em, item-indent=1.8em, column-sep=0.6em, after-item-skip=5pt, before-skip=3pt, after-skip=4pt}

% ---------- capçalera i peu ----------
\newcommand{\llicencia}{\href{https://creativecommons.org/licenses/by-nc-sa/4.0/deed.es}{\ccbyncsa}\ \ Licencia Creative Commons Atribución-NoComercial-CompartirIgual 4.0 Internacional (CC BY-NC-SA 4.0)}
\pagestyle{fancy}
\fancyhf{}
\fancyhead[L]{\raisebox{-0.35\height}{\includegraphics[height=0.95cm]{logo-gva.png}}}
\fancyhead[C]{\small\color{muted}Ficha de ejercicios · Unidad 7 · 1\textsuperscript{o} ESO}
\fancyhead[R]{\raisebox{-0.35\height}{\includegraphics[height=0.95cm]{logo-vermellar.png}}}
\fancyfoot[L]{\scriptsize\color{muted}\llicencia}
\fancyfoot[R]{\small\color{muted}\thepage}
\renewcommand{\headrulewidth}{0.4pt}
\renewcommand{\headrule}{\color{muted!40}\hrule height \headrulewidth}
\fancypagestyle{primera}{\fancyhf{}\renewcommand{\headrulewidth}{0pt}
  \fancyfoot[L]{\scriptsize\color{muted}\llicencia}
  \fancyfoot[R]{\small\color{muted}\thepage}}

% ---------- nivells i exercicis ----------
\newcommand{\estrelles}[1]{{\color{accent}\ifcase#1\or\faStar\faStar[regular]\faStar[regular]\or\faStar\faStar\faStar[regular]\or\faStar\faStar\faStar\fi}}
\newcounter{exer}
\newcommand{\ex}[1]{\par\needspace{4\baselineskip}\medskip\stepcounter{exer}\noindent{\color{navy}\bfseries\theexer.}\ \ {\small\estrelles{#1}}\ \ }

\newcommand{\seccio}[2]{\par\needspace{10\baselineskip}\bigskip
  \begin{tcolorbox}[colback=navy,colframe=navy,coltext=white,arc=2mm,boxrule=0pt,left=3mm,right=3mm,top=1.2mm,bottom=1.2mm]
  {\large\bfseries #1\quad #2}
  \end{tcolorbox}\vspace{-2pt}}

\newcolumntype{C}[1]{>{\centering\arraybackslash}p{#1}}

% quadre de recordatori (només a la fitxa d'exercicis)
\newcommand{\recorda}[1]{\par\smallskip\begin{tcolorbox}[colback=accent!7,colframe=accent!55,arc=2mm,boxrule=0.5pt,left=2.5mm,right=2.5mm,top=1mm,bottom=1mm]\small{\color{accent}\faLightbulb[regular]}\ \textbf{Recuerda:}\ #1\end{tcolorbox}}
\newcommand{\res}[1]{#1}\newcommand{\rest}[1]{#1}

\def\fmneed{14}
% =====================================================================
% Capçalera FentMates de les fitxes: un color per a cada curs.
% Els colors són els de la pàgina inicial (LEVELS a assets/js/materials.js);
% el logo ve de _eines/logo/pdf/ (python3 _eines/logo/genera.py).
%
% Ús, al final del preàmbul de la unitat:
%   % =====================================================================
% Capçalera FentMates de les fitxes: un color per a cada curs.
% Els colors són els de la pàgina inicial (LEVELS a assets/js/materials.js);
% el logo ve de _eines/logo/pdf/ (python3 _eines/logo/genera.py).
%
% Ús, al final del preàmbul de la unitat:
%   % =====================================================================
% Capçalera FentMates de les fitxes: un color per a cada curs.
% Els colors són els de la pàgina inicial (LEVELS a assets/js/materials.js);
% el logo ve de _eines/logo/pdf/ (python3 _eines/logo/genera.py).
%
% Ús, al final del preàmbul de la unitat:
%   \input{../comu/fentmates.tex}
%   \fmcurs{1eso}
% i, al començament del document:
%   \capcalera{Matemàtiques}{1r ESO}{Fitxa d'exercicis · Unitat 2}{Divisibilitat}
% Opcional abans de \input: \def\fmneed{14} (espai mínim abans de cada secció).
% =====================================================================
\makeatletter
\def\fm@logodir{../../logo/pdf}
\providecommand{\fmneed}{10}
\def\fm@cursnom{}

% clau del curs, color i color del text damunt del color
\newcommand{\fm@curs}[3]{\@namedef{fm@color@#1}{#2}\@namedef{fm@text@#1}{#3}}
\fm@curs{1eso}{E6B84A}{navy}
\fm@curs{2eso}{3B82F6}{white}
\fm@curs{3eso}{FF7A50}{white}
\fm@curs{4eso-a}{A78BFA}{white}
\fm@curs{4eso-b}{6D28D9}{white}
\fm@curs{1batx-mi}{9B1C1C}{white}
\fm@curs{1batx-ccss}{EF6B6B}{white}
\fm@curs{2batx-mii}{1E6B47}{white}
\fm@curs{2batx-ccss}{2BB3A3}{white}

\newcommand{\fmcurs}[1]{%
  \@ifundefined{fm@color@#1}{\PackageError{fentmates}{Curs desconegut: #1}{}}{}%
  \edef\fm@tmp{\noexpand\definecolor{curs}{HTML}{\@nameuse{fm@color@#1}}%
               \noexpand\colorlet{curstext}{\@nameuse{fm@text@#1}}}\fm@tmp
  \colorlet{cursfosc}{curs!72!black}%
  \def\fm@logo{\fm@logodir/vertical-#1.pdf}%
  \def\fm@marca{\fm@logodir/marca-#1.pdf}}

% ---------- capçalera de la primera pàgina ----------
% #1 matèria (petit), #2 curs, #3 tipus de document i unitat, #4 títol
\newcommand{\capcalera}[4]{%
  \gdef\fm@cursnom{#2}%
  \par\noindent
  \begin{tcolorbox}[colback=curs,colframe=curs,boxrule=0pt,arc=2.2mm,width=3.5cm,height=2.15cm,
      valign=center,left=3mm,right=2mm,top=1mm,bottom=1mm,nobeforeafter,box align=center]
    \color{curstext}{\footnotesize #1}\par\vspace{1pt}{\fontsize{19}{21}\selectfont\bfseries #2\par}%
  \end{tcolorbox}\hspace{5mm}%
  \begin{minipage}[c]{\dimexpr\linewidth-3.5cm-5mm-2.9cm\relax}
    {\color{cursfosc}\small\bfseries\MakeUppercase{#3}}\par\vspace{2pt}
    {\color{navy}\fontsize{20}{23}\selectfont\bfseries #4\par}
  \end{minipage}\hfill
  \begin{minipage}[c]{2.6cm}\raggedleft\includegraphics[width=2.6cm]{\fm@logo}\end{minipage}\par
  \vspace{2pt}}

% ---------- barres de secció del color del curs ----------
\renewcommand{\seccio}[2]{\par\needspace{\fmneed\baselineskip}\bigskip
  \begin{tcolorbox}[colback=curs,colframe=curs,coltext=curstext,arc=2mm,boxrule=0pt,left=3mm,right=3mm,top=1.2mm,bottom=1.2mm]
  {\large\bfseries #1\quad #2}
  \end{tcolorbox}\vspace{-2pt}}

% ---------- estrelles de dificultat del color del curs (a Batxillerat no se'n posen) ----------
\@ifundefined{estrelles}{}{%
  \renewcommand{\estrelles}[1]{{\color{curs}\ifcase#1\or\faStar\faStar[regular]\faStar[regular]\or\faStar\faStar\faStar[regular]\or\faStar\faStar\faStar\fi}}}

% ---------- capçalera i peu de les pàgines ----------
% primera pàgina: logos institucionals al peu
\fancypagestyle{primera}{\fancyhf{}\renewcommand{\headrulewidth}{0pt}
  \fancyfoot[L]{\raisebox{-0.32\height}{\includegraphics[height=0.78cm]{logo-gva.png}}\hspace{5mm}%
    \raisebox{-0.32\height}{\includegraphics[height=0.78cm]{logo-vermellar.png}}}
  \fancyfoot[R]{\scriptsize\color{muted}\href{https://creativecommons.org/licenses/by-nc-sa/4.0/deed.ca}{\ccbyncsa}\ \ CC BY-NC-SA 4.0 · fentmates.org\qquad\small\thepage}}
% la resta: el símbol a l'esquerra i el curs a la dreta, amb el títol al centre com abans
\fancyhead[L]{\raisebox{-0.3\height}{\includegraphics[height=0.62cm]{\fm@marca}}}
\fancyhead[R]{\small\bfseries\color{cursfosc}\fm@cursnom}
\renewcommand{\headrule}{\color{curs!45}\hrule height \headrulewidth}
\makeatother

%   \fmcurs{1eso}
% i, al començament del document:
%   \capcalera{Matemàtiques}{1r ESO}{Fitxa d'exercicis · Unitat 2}{Divisibilitat}
% Opcional abans de \input: \def\fmneed{14} (espai mínim abans de cada secció).
% =====================================================================
\makeatletter
\def\fm@logodir{../../logo/pdf}
\providecommand{\fmneed}{10}
\def\fm@cursnom{}

% clau del curs, color i color del text damunt del color
\newcommand{\fm@curs}[3]{\@namedef{fm@color@#1}{#2}\@namedef{fm@text@#1}{#3}}
\fm@curs{1eso}{E6B84A}{navy}
\fm@curs{2eso}{3B82F6}{white}
\fm@curs{3eso}{FF7A50}{white}
\fm@curs{4eso-a}{A78BFA}{white}
\fm@curs{4eso-b}{6D28D9}{white}
\fm@curs{1batx-mi}{9B1C1C}{white}
\fm@curs{1batx-ccss}{EF6B6B}{white}
\fm@curs{2batx-mii}{1E6B47}{white}
\fm@curs{2batx-ccss}{2BB3A3}{white}

\newcommand{\fmcurs}[1]{%
  \@ifundefined{fm@color@#1}{\PackageError{fentmates}{Curs desconegut: #1}{}}{}%
  \edef\fm@tmp{\noexpand\definecolor{curs}{HTML}{\@nameuse{fm@color@#1}}%
               \noexpand\colorlet{curstext}{\@nameuse{fm@text@#1}}}\fm@tmp
  \colorlet{cursfosc}{curs!72!black}%
  \def\fm@logo{\fm@logodir/vertical-#1.pdf}%
  \def\fm@marca{\fm@logodir/marca-#1.pdf}}

% ---------- capçalera de la primera pàgina ----------
% #1 matèria (petit), #2 curs, #3 tipus de document i unitat, #4 títol
\newcommand{\capcalera}[4]{%
  \gdef\fm@cursnom{#2}%
  \par\noindent
  \begin{tcolorbox}[colback=curs,colframe=curs,boxrule=0pt,arc=2.2mm,width=3.5cm,height=2.15cm,
      valign=center,left=3mm,right=2mm,top=1mm,bottom=1mm,nobeforeafter,box align=center]
    \color{curstext}{\footnotesize #1}\par\vspace{1pt}{\fontsize{19}{21}\selectfont\bfseries #2\par}%
  \end{tcolorbox}\hspace{5mm}%
  \begin{minipage}[c]{\dimexpr\linewidth-3.5cm-5mm-2.9cm\relax}
    {\color{cursfosc}\small\bfseries\MakeUppercase{#3}}\par\vspace{2pt}
    {\color{navy}\fontsize{20}{23}\selectfont\bfseries #4\par}
  \end{minipage}\hfill
  \begin{minipage}[c]{2.6cm}\raggedleft\includegraphics[width=2.6cm]{\fm@logo}\end{minipage}\par
  \vspace{2pt}}

% ---------- barres de secció del color del curs ----------
\renewcommand{\seccio}[2]{\par\needspace{\fmneed\baselineskip}\bigskip
  \begin{tcolorbox}[colback=curs,colframe=curs,coltext=curstext,arc=2mm,boxrule=0pt,left=3mm,right=3mm,top=1.2mm,bottom=1.2mm]
  {\large\bfseries #1\quad #2}
  \end{tcolorbox}\vspace{-2pt}}

% ---------- estrelles de dificultat del color del curs (a Batxillerat no se'n posen) ----------
\@ifundefined{estrelles}{}{%
  \renewcommand{\estrelles}[1]{{\color{curs}\ifcase#1\or\faStar\faStar[regular]\faStar[regular]\or\faStar\faStar\faStar[regular]\or\faStar\faStar\faStar\fi}}}

% ---------- capçalera i peu de les pàgines ----------
% primera pàgina: logos institucionals al peu
\fancypagestyle{primera}{\fancyhf{}\renewcommand{\headrulewidth}{0pt}
  \fancyfoot[L]{\raisebox{-0.32\height}{\includegraphics[height=0.78cm]{logo-gva.png}}\hspace{5mm}%
    \raisebox{-0.32\height}{\includegraphics[height=0.78cm]{logo-vermellar.png}}}
  \fancyfoot[R]{\scriptsize\color{muted}\href{https://creativecommons.org/licenses/by-nc-sa/4.0/deed.ca}{\ccbyncsa}\ \ CC BY-NC-SA 4.0 · fentmates.org\qquad\small\thepage}}
% la resta: el símbol a l'esquerra i el curs a la dreta, amb el títol al centre com abans
\fancyhead[L]{\raisebox{-0.3\height}{\includegraphics[height=0.62cm]{\fm@marca}}}
\fancyhead[R]{\small\bfseries\color{cursfosc}\fm@cursnom}
\renewcommand{\headrule}{\color{curs!45}\hrule height \headrulewidth}
\makeatother

%   \fmcurs{1eso}
% i, al començament del document:
%   \capcalera{Matemàtiques}{1r ESO}{Fitxa d'exercicis · Unitat 2}{Divisibilitat}
% Opcional abans de \input: \def\fmneed{14} (espai mínim abans de cada secció).
% =====================================================================
\makeatletter
\def\fm@logodir{../../logo/pdf}
\providecommand{\fmneed}{10}
\def\fm@cursnom{}

% clau del curs, color i color del text damunt del color
\newcommand{\fm@curs}[3]{\@namedef{fm@color@#1}{#2}\@namedef{fm@text@#1}{#3}}
\fm@curs{1eso}{E6B84A}{navy}
\fm@curs{2eso}{3B82F6}{white}
\fm@curs{3eso}{FF7A50}{white}
\fm@curs{4eso-a}{A78BFA}{white}
\fm@curs{4eso-b}{6D28D9}{white}
\fm@curs{1batx-mi}{9B1C1C}{white}
\fm@curs{1batx-ccss}{EF6B6B}{white}
\fm@curs{2batx-mii}{1E6B47}{white}
\fm@curs{2batx-ccss}{2BB3A3}{white}

\newcommand{\fmcurs}[1]{%
  \@ifundefined{fm@color@#1}{\PackageError{fentmates}{Curs desconegut: #1}{}}{}%
  \edef\fm@tmp{\noexpand\definecolor{curs}{HTML}{\@nameuse{fm@color@#1}}%
               \noexpand\colorlet{curstext}{\@nameuse{fm@text@#1}}}\fm@tmp
  \colorlet{cursfosc}{curs!72!black}%
  \def\fm@logo{\fm@logodir/vertical-#1.pdf}%
  \def\fm@marca{\fm@logodir/marca-#1.pdf}}

% ---------- capçalera de la primera pàgina ----------
% #1 matèria (petit), #2 curs, #3 tipus de document i unitat, #4 títol
\newcommand{\capcalera}[4]{%
  \gdef\fm@cursnom{#2}%
  \par\noindent
  \begin{tcolorbox}[colback=curs,colframe=curs,boxrule=0pt,arc=2.2mm,width=3.5cm,height=2.15cm,
      valign=center,left=3mm,right=2mm,top=1mm,bottom=1mm,nobeforeafter,box align=center]
    \color{curstext}{\footnotesize #1}\par\vspace{1pt}{\fontsize{19}{21}\selectfont\bfseries #2\par}%
  \end{tcolorbox}\hspace{5mm}%
  \begin{minipage}[c]{\dimexpr\linewidth-3.5cm-5mm-2.9cm\relax}
    {\color{cursfosc}\small\bfseries\MakeUppercase{#3}}\par\vspace{2pt}
    {\color{navy}\fontsize{20}{23}\selectfont\bfseries #4\par}
  \end{minipage}\hfill
  \begin{minipage}[c]{2.6cm}\raggedleft\includegraphics[width=2.6cm]{\fm@logo}\end{minipage}\par
  \vspace{2pt}}

% ---------- barres de secció del color del curs ----------
\renewcommand{\seccio}[2]{\par\needspace{\fmneed\baselineskip}\bigskip
  \begin{tcolorbox}[colback=curs,colframe=curs,coltext=curstext,arc=2mm,boxrule=0pt,left=3mm,right=3mm,top=1.2mm,bottom=1.2mm]
  {\large\bfseries #1\quad #2}
  \end{tcolorbox}\vspace{-2pt}}

% ---------- estrelles de dificultat del color del curs (a Batxillerat no se'n posen) ----------
\@ifundefined{estrelles}{}{%
  \renewcommand{\estrelles}[1]{{\color{curs}\ifcase#1\or\faStar\faStar[regular]\faStar[regular]\or\faStar\faStar\faStar[regular]\or\faStar\faStar\faStar\fi}}}

% ---------- capçalera i peu de les pàgines ----------
% primera pàgina: logos institucionals al peu
\fancypagestyle{primera}{\fancyhf{}\renewcommand{\headrulewidth}{0pt}
  \fancyfoot[L]{\raisebox{-0.32\height}{\includegraphics[height=0.78cm]{logo-gva.png}}\hspace{5mm}%
    \raisebox{-0.32\height}{\includegraphics[height=0.78cm]{logo-vermellar.png}}}
  \fancyfoot[R]{\scriptsize\color{muted}\href{https://creativecommons.org/licenses/by-nc-sa/4.0/deed.ca}{\ccbyncsa}\ \ CC BY-NC-SA 4.0 · fentmates.org\qquad\small\thepage}}
% la resta: el símbol a l'esquerra i el curs a la dreta, amb el títol al centre com abans
\fancyhead[L]{\raisebox{-0.3\height}{\includegraphics[height=0.62cm]{\fm@marca}}}
\fancyhead[R]{\small\bfseries\color{cursfosc}\fm@cursnom}
\renewcommand{\headrule}{\color{curs!45}\hrule height \headrulewidth}
\makeatother

\fmcurs{1eso}

\begin{document}
\thispagestyle{primera}

\capcalera{Matemáticas}{1\textsuperscript{o} ESO}{Ficha de ejercicios · Unidad 7}{Geometría básica del plano}
\vspace{10pt}
Nombre y apellidos: \rule{6.6cm}{0.4pt}\hfill Grupo: \rule{1.6cm}{0.4pt}\hfill Fecha: \rule{2.2cm}{0.4pt}
\vspace{8pt}

\begin{tcolorbox}[colback=soft,colframe=soft,arc=2mm,left=3mm,right=3mm,top=1.5mm,bottom=1.5mm]
\small Ejercicios graduados por niveles de dificultad:\quad \estrelles{1}\ básico \qquad \estrelles{2}\ medio \qquad \estrelles{3}\ avanzado.\\[2pt] Resuélvelos en el cuaderno mostrando todos los pasos. Ten a mano la regla, el compás y el transportador de ángulos, y haz un dibujo siempre que te ayude.
\end{tcolorbox}

% =====================================================
\seccio{1}{Puntos, rectas, semirrectas y segmentos}

\ex{1} Observa la figura y responde:\par\smallskip\begin{center}\begin{tikzpicture}[x=0.6cm,y=0.6cm,baseline=(current bounding box.center),line join=round,line cap=round]
\draw[thick,latex-latex] (-0.500,0.000) -- (7.500,0.000);
\node[font=\small,below] at (7.300,0.000) {$r$};
\fill (1.000,0.000) circle (1.5pt);
\node[font=\small,above] at (1.000,0.000) {$A$};
\fill (3.500,0.000) circle (1.5pt);
\node[font=\small,above] at (3.500,0.000) {$B$};
\fill (6.000,0.000) circle (1.5pt);
\node[font=\small,above] at (6.000,0.000) {$C$};
\fill (4.600,1.500) circle (1.5pt);
\node[font=\small,above] at (4.600,1.500) {$D$};
\end{tikzpicture}\end{center}\begin{tasks}(1)
  \task ¿Qué puntos son de la recta $r$? ¿Qué punto no está en ella?
  \task Escribe todos los segmentos que tienen los extremos en dos de los puntos $A$, $B$ y $C$.
  \task ¿Cuántas semirrectas de la recta $r$ tienen el origen en $B$?
\end{tasks}\recorda{La \textbf{recta} no tiene principio ni final; la \textbf{semirrecta} tiene origen pero no tiene final; el \textbf{segmento} tiene dos extremos y se puede medir. Por dos puntos pasa una única recta.}

\ex{1} Di si cada figura es una recta, una semirrecta o un segmento:\begin{tasks}(3)
  \task \begin{tikzpicture}[x=0.6cm,y=0.6cm,baseline=(current bounding box.center),line join=round,line cap=round]
\draw[thick,latex-latex] (0.000,0.000) -- (3.000,0.600);
\fill (1.000,0.200) circle (1.5pt);
\fill (2.000,0.400) circle (1.5pt);
\end{tikzpicture}
  \task \begin{tikzpicture}[x=0.6cm,y=0.6cm,baseline=(current bounding box.center),line join=round,line cap=round]
\draw[thick,-] (0.000,0.000) -- (3.000,0.600);
\fill (0.000,0.000) circle (1.5pt);
\fill (3.000,0.600) circle (1.5pt);
\end{tikzpicture}
  \task \begin{tikzpicture}[x=0.6cm,y=0.6cm,baseline=(current bounding box.center),line join=round,line cap=round]
\draw[thick,-latex] (0.000,0.000) -- (3.000,0.600);
\fill (0.000,0.000) circle (1.5pt);
\end{tikzpicture}
\end{tasks}

\ex{1} Responde:\begin{tasks}(1)
  \task El segmento $\overline{AB}$ mide 7\,cm y $M$ es su punto medio. ¿Cuánto mide $\overline{AM}$?
  \task $M$ es el punto medio de $\overline{PQ}$ y $\overline{PM}$ mide 4,2\,cm. ¿Cuánto mide $\overline{PQ}$?
\end{tasks}

\ex{2} Sobre una recta marcamos cuatro puntos. ¿Cuántos segmentos determinan? ¿Y si marcamos seis?

\ex{2} Los puntos $A$, $B$ y $C$ están sobre una recta, en este orden. $\overline{AB}=5$\,cm y $\overline{BC}=3$\,cm. ¿Cuánto mide $\overline{AC}$? ¿Y si $C$ estuviera entre $A$ y $B$ (con $\overline{BC}=3$\,cm)?

\ex{3} Dibujamos puntos de manera que no haya tres alineados. ¿Cuántas rectas podemos trazar que pasen por dos de ellos si hay 4 puntos? ¿Y 5? ¿Y 7?


% =====================================================
\seccio{2}{Posiciones relativas de rectas. Mediatriz}

\ex{1} Las rectas $r$ y $s$ son paralelas. Di cómo son:\par\smallskip\begin{center}\begin{tikzpicture}[x=0.6cm,y=0.6cm,baseline=(current bounding box.center),line join=round,line cap=round]
\draw[thick,latex-latex] (-0.400,0.000) -- (6.400,0.000);
\draw[thick,latex-latex] (-0.400,1.600) -- (6.400,1.600);
\draw[thick,latex-latex] (1.500,-0.900) -- (1.500,2.600);
\draw[thick,latex-latex] (2.800,-0.900) -- (6.200,2.600);
\draw[thin] (1.867,0.000) -- (1.867,0.367) -- (1.500,0.367);
\node[font=\small,right] at (6.400,0.000) {$r$};
\node[font=\small,right] at (6.400,1.600) {$s$};
\node[font=\small,above] at (1.500,2.600) {$t$};
\node[font=\small,above] at (6.200,2.600) {$u$};
\end{tikzpicture}\end{center}\begin{tasks}(4)
  \task $r$ y $s$
  \task $r$ y $t$
  \task $s$ y $t$
  \task $r$ y $u$
\end{tasks}\recorda{Dos rectas son \textbf{secantes} si se cortan en un punto, \textbf{paralelas} si no se cortan nunca y \textbf{perpendiculares} si se cortan formando cuatro ángulos rectos ($r\perp s$).}

\ex{1} Di si son paralelas, secantes o perpendiculares:\begin{tasks}(2)
  \task Dos lados opuestos de una hoja de papel.
  \task Dos lados de una hoja que forman una esquina.
  \task Las agujas de un reloj a las 9 en punto.
  \task Las agujas de un reloj a las 10 en punto.
\end{tasks}

\ex{1} Dibuja un segmento $\overline{AB}$ de 6\,cm y traza su mediatriz con regla y compás. ¿A qué distancia de $A$ corta el segmento?

\ex{2} Di si es cierto o falso:\begin{tasks}(1)
  \task Dos rectas perpendiculares son siempre secantes.
  \task Dos rectas secantes son siempre perpendiculares.
  \task Si $r\parallel s$ y $s\parallel t$, entonces $r\parallel t$.
  \task Si $r\perp s$ y $s\perp t$, entonces $r\perp t$.
\end{tasks}

\ex{2} El punto $P$ está sobre la mediatriz del segmento $\overline{AB}$, que mide 8\,cm, y la distancia de $P$ a $A$ es 5\,cm. ¿Cuánto mide $\overline{PB}$? ¿A qué distancia de $A$ está el punto medio $M$?

\ex{3} Dos casas, $A$ y $B$, están al mismo lado de un camino recto. ¿Dónde se tiene que poner una fuente, junto al camino, para que esté a la misma distancia de las dos casas? Explica cómo la encontrarías.


% =====================================================
\seccio{3}{Ángulos: elementos, medida y clasificación}

\ex{1} Clasifica estos ángulos (nulo, agudo, recto, obtuso, llano, cóncavo o completo):\begin{tasks}(4)
  \task $25^\circ$
  \task $90^\circ$
  \task $100^\circ$
  \task $180^\circ$
  \task $270^\circ$
  \task $360^\circ$
  \task $179^\circ$
  \task $89^\circ$
\end{tasks}\recorda{Agudo: menos de $90^\circ$. Recto: $90^\circ$. Obtuso: entre $90^\circ$ y $180^\circ$. Llano: $180^\circ$. Cóncavo: más de $180^\circ$. Completo: $360^\circ$.}

\ex{1} Clasifica cada ángulo y mídelo con el transportador:\begin{tasks}(2)
  \task \begin{tikzpicture}[x=0.8cm,y=0.8cm,baseline=(current bounding box.center),line join=round,line cap=round]
\draw[fill=accent!30,draw=accent] (0.000,0.000) -- (0.591,0.104) arc[start angle=10,end angle=50,radius=0.6] -- cycle;
\draw[thick,-latex] (0.000,0.000) -- (2.954,0.521);
\draw[thick,-latex] (0.000,0.000) -- (1.928,2.298);
\fill (0.000,0.000) circle (1.5pt);
\end{tikzpicture}
  \task \begin{tikzpicture}[x=0.8cm,y=0.8cm,baseline=(current bounding box.center),line join=round,line cap=round]
\draw[fill=accent!30,draw=accent] (0.000,0.000) -- (0.400,0.000) -- (0.400,0.400) -- (0.000,0.400) -- cycle;
\draw[thick,-latex] (0.000,0.000) -- (3.000,0.000);
\draw[thick,-latex] (0.000,0.000) -- (0.000,3.000);
\fill (0.000,0.000) circle (1.5pt);
\end{tikzpicture}
  \task \begin{tikzpicture}[x=0.8cm,y=0.8cm,baseline=(current bounding box.center),line join=round,line cap=round]
\draw[fill=accent!30,draw=accent] (0.000,0.000) -- (0.600,0.000) arc[start angle=0,end angle=140,radius=0.6] -- cycle;
\draw[thick,-latex] (0.000,0.000) -- (3.000,0.000);
\draw[thick,-latex] (0.000,0.000) -- (-2.298,1.928);
\fill (0.000,0.000) circle (1.5pt);
\end{tikzpicture}
  \task \begin{tikzpicture}[x=0.8cm,y=0.8cm,baseline=(current bounding box.center),line join=round,line cap=round]
\draw[fill=accent!30,draw=accent] (0.000,0.000) -- (0.600,0.000) arc[start angle=0,end angle=230,radius=0.6] -- cycle;
\draw[thick,-latex] (0.000,0.000) -- (3.000,0.000);
\draw[thick,-latex] (0.000,0.000) -- (-1.928,-2.298);
\fill (0.000,0.000) circle (1.5pt);
\end{tikzpicture}
\end{tasks}

\ex{1} ¿Qué ángulo (convexo) forman las agujas de un reloj a las horas en punto siguientes?\begin{tasks}(6)
  \task 3 h
  \task 4 h
  \task 5 h
  \task 6 h
  \task 9 h
  \task 11 h
\end{tasks}

\ex{2} ¿Cuántos grados gira la aguja de los minutos en estos tiempos?\begin{tasks}(4)
  \task 10 min
  \task 25 min
  \task 1 hora
  \task tres cuartos de hora
\end{tasks}

\ex{2} Responde:\begin{tasks}(1)
  \task ¿Qué ángulos forma la bisectriz de un ángulo de $136^\circ$ con cada uno de los lados?
  \task La bisectriz de un ángulo forma con uno de los lados un ángulo de $43^\circ$. ¿Cuánto mide el ángulo?
\end{tasks}

\ex{2} Dos semirrectas con el mismo origen forman un ángulo convexo de $125^\circ$. ¿Cuánto mide el ángulo cóncavo que forman?

\ex{3} Responde:\begin{tasks}(1)
  \task Dividimos un ángulo completo en 5 ángulos iguales. ¿Cuánto mide cada uno? ¿Y si lo dividimos en 8?
  \task ¿Cuántas veces cabe un ángulo de $15^\circ$ en un ángulo llano?
\end{tasks}


% =====================================================
\seccio{4}{Medida de ángulos: el sistema sexagesimal}

\ex{1} Expresa en segundos:\begin{tasks}(3)
  \task $2^\circ$
  \task $15'$
  \task $1^\circ\,20'$
  \task $3^\circ\,5'\,40''$
  \task $10^\circ\,10'\,10''$
  \task $45'\,30''$
\end{tasks}\recorda{$1^\circ=60'$, \ $1'=60''$, \ $1^\circ=3600''$. Para pasar a una unidad más pequeña multiplicamos; para pasar a una unidad más grande dividimos entre 60.}

\ex{1} Expresa en minutos:\begin{tasks}(4)
  \task $5^\circ$
  \task $2^\circ\,30'$
  \task $12^\circ\,45'$
  \task $180^\circ$
\end{tasks}

\ex{1} Pasa a forma compleja (grados, minutos y segundos):\begin{tasks}(3)
  \task $150'$
  \task $4000''$
  \task $7385''$
  \task $10\,000''$
  \task $1000'$
  \task $86\,400''$
\end{tasks}

\ex{2} Expresa en grados y minutos:\begin{tasks}(4)
  \task $12,5^\circ$
  \task $30,25^\circ$
  \task $8,75^\circ$
  \task $0,1^\circ$
\end{tasks}

\ex{2} Calcula:\begin{tasks}(3)
  \task $23^\circ\,45'+18^\circ\,30'$
  \task $45^\circ\,36'\,50''+27^\circ\,48'\,25''$
  \task $120^\circ\,15'\,40''+59^\circ\,44'\,20''$
\end{tasks}

\ex{2} Calcula:\begin{tasks}(3)
  \task $75^\circ\,40'-32^\circ\,15'$
  \task $60^\circ\,20'\,15''-25^\circ\,45'\,30''$
  \task $180^\circ-97^\circ\,34'\,18''$
\end{tasks}

\ex{2} Calcula:\begin{tasks}(3)
  \task $(12^\circ\,20')\cdot3$
  \task $(25^\circ\,15'\,40'')\cdot4$
  \task $(8^\circ\,45'\,50'')\cdot6$
\end{tasks}

\ex{3} Calcula y comprueba el resultado multiplicando:\begin{tasks}(3)
  \task $(75^\circ\,36'):4$
  \task $(100^\circ\,20'\,30''):5$
  \task $(47^\circ\,13'):3$
\end{tasks}

\ex{3} La Tierra da una vuelta completa sobre sí misma en 24 horas. ¿Cuántos grados gira en 1 hora? ¿Y en 1 minuto? ¿Y en 1 segundo?


% =====================================================
\seccio{5}{Relaciones entre ángulos}

\ex{1} Calcula el complementario de cada ángulo:\begin{tasks}(4)
  \task $30^\circ$
  \task $47^\circ$
  \task $89^\circ$
  \task $15^\circ\,30'$
\end{tasks}\recorda{Dos ángulos son \textbf{complementarios} si suman $90^\circ$ y \textbf{suplementarios} si suman $180^\circ$. Los ángulos \textbf{opuestos por el vértice} son iguales y los \textbf{adyacentes} suman $180^\circ$.}

\ex{1} Calcula el suplementario de cada ángulo:\begin{tasks}(4)
  \task $45^\circ$
  \task $120^\circ$
  \task $179^\circ$
  \task $92^\circ\,15'\,40''$
\end{tasks}

\ex{1} Dos rectas se cortan y uno de los ángulos mide $35^\circ$. Calcula los ángulos $a$, $b$ y $c$ y explica cómo los has encontrado.\par\smallskip\begin{center}\begin{tikzpicture}[x=0.6cm,y=0.6cm,baseline=(current bounding box.center),line join=round,line cap=round]
\draw[fill=accent!30,draw=accent] (0.000,0.000) -- (0.433,0.250) arc[start angle=30,end angle=65,radius=0.5] -- cycle;
\draw[thick,latex-latex] (-1.732,-1.000) -- (1.732,1.000);
\draw[thick,latex-latex] (-0.845,-1.813) -- (0.845,1.813);
\fill (0.000,0.000) circle (1.5pt);
\node[font=\small,inner sep=1pt] at (0.648,0.695) {$35^\circ$};
\node[font=\small,inner sep=1pt] at (-0.629,0.408) {$a$};
\node[font=\small,inner sep=1pt] at (-0.511,-0.549) {$b$};
\node[font=\small,inner sep=1pt] at (0.629,-0.408) {$c$};
\end{tikzpicture}\end{center}

\ex{2} Las rectas $r$ y $s$ son paralelas y el ángulo 1 mide $58^\circ$. Calcula el resto de ángulos numerados.\par\smallskip\begin{center}\begin{tikzpicture}[x=0.85cm,y=0.85cm,baseline=(current bounding box.center),line join=round,line cap=round]
\draw[thick,latex-latex] (-2.200,0.000) -- (3.325,0.000);
\draw[thick,latex-latex] (-2.200,1.800) -- (3.325,1.800);
\draw[thick,latex-latex] (-0.689,-1.102) -- (1.814,2.902);
\node[font=\small,above] at (-2.200,1.800) {$r$};
\node[font=\small,above] at (-2.200,0.000) {$s$};
\node[font=\small,inner sep=1pt] at (0.542,0.301) {\scriptsize 1};
\node[font=\small,inner sep=1pt] at (-0.301,0.542) {\scriptsize 2};
\node[font=\small,inner sep=1pt] at (-0.542,-0.301) {\scriptsize 3};
\node[font=\small,inner sep=1pt] at (0.301,-0.542) {\scriptsize 4};
\node[font=\small,inner sep=1pt] at (1.667,2.101) {\scriptsize 5};
\node[font=\small,inner sep=1pt] at (0.824,2.342) {\scriptsize 6};
\node[font=\small,inner sep=1pt] at (0.583,1.499) {\scriptsize 7};
\node[font=\small,inner sep=1pt] at (1.425,1.258) {\scriptsize 8};
\end{tikzpicture}\end{center}

\ex{2} Di si es cierto o falso:\begin{tasks}(1)
  \task Dos ángulos agudos pueden ser suplementarios.
  \task El suplementario de un ángulo agudo es obtuso.
  \task Dos ángulos adyacentes siempre son suplementarios.
  \task Un ángulo obtuso tiene complementario.
\end{tasks}

\ex{2} Calcula el complementario y el suplementario del ángulo $37^\circ\,25'\,48''$.

\ex{3} Dos ángulos son complementarios y uno de ellos mide $20^\circ$ más que el otro. ¿Cuánto mide cada uno?

\ex{3} El suplementario de un ángulo es el cuádruple del ángulo. ¿Cuánto mide el ángulo?

\ex{3} Dos rectas paralelas cortadas por una secante forman dos ángulos alternos internos que miden $3x+10^\circ$ y $70^\circ$. ¿Cuánto vale $x$?


% =====================================================
\seccio{6}{Polígonos: elementos y clasificación}

\ex{1} Escribe el nombre del polígono que tiene:\begin{tasks}(3)
  \task 5 lados
  \task 8 lados
  \task 6 lados
  \task 10 lados
  \task 7 lados
  \task 12 lados
\end{tasks}

\ex{1} Para cada polígono, di cuántos lados tiene, si es convexo o cóncavo y si es regular o irregular:\begin{tasks}(4)
  \task \begin{tikzpicture}[x=0.6cm,y=0.6cm,baseline=(current bounding box.center),line join=round,line cap=round]
\draw[fill=teal!12,draw=teal,thick] (-0.647,-0.890) -- (0.647,-0.890) -- (1.046,0.340) -- (0.000,1.100) -- (-1.046,0.340) -- cycle;
\end{tikzpicture}
  \task \begin{tikzpicture}[x=0.6cm,y=0.6cm,baseline=(current bounding box.center),line join=round,line cap=round]
\draw[fill=teal!12,draw=teal,thick] (0.000,0.000) -- (2.400,0.000) -- (2.400,2.000) -- (1.200,0.900) -- (0.000,2.000) -- cycle;
\end{tikzpicture}
  \task \begin{tikzpicture}[x=0.6cm,y=0.6cm,baseline=(current bounding box.center),line join=round,line cap=round]
\draw[fill=teal!12,draw=teal,thick] (0.000,0.000) -- (2.800,0.000) -- (2.800,1.600) -- (0.000,1.600) -- cycle;
\end{tikzpicture}
  \task \begin{tikzpicture}[x=0.6cm,y=0.6cm,baseline=(current bounding box.center),line join=round,line cap=round]
\draw[fill=teal!12,draw=teal,thick] (0.000,0.000) -- (2.400,0.300) -- (2.000,1.900) -- (0.300,1.500) -- cycle;
\end{tikzpicture}
\end{tasks}\recorda{Un polígono es \textbf{cóncavo} si tiene algún ángulo interior de más de $180^\circ$, y \textbf{regular} si tiene todos los lados \emph{y} todos los ángulos iguales.}

\ex{1} ¿Cuántas diagonales salen de cada vértice y cuántas tiene en total?\begin{tasks}(4)
  \task cuadrilátero
  \task pentágono
  \task octógono
  \task decágono
\end{tasks}

\ex{2} Di si es cierto o falso:\begin{tasks}(1)
  \task Un polígono regular tiene todos los lados iguales.
  \task Un rombo (4 lados iguales) es un polígono regular.
  \task Un polígono cóncavo tiene algún ángulo interior de más de $180^\circ$.
  \task Un triángulo tiene 3 diagonales.
\end{tasks}

\ex{2} ¿Qué polígono tiene 14 diagonales? ¿Y cuál tiene 27?

\ex{3} De cada vértice de un polígono salen 9 diagonales. ¿Cuántos lados tiene? ¿Cuántas diagonales tiene en total?


% =====================================================
\seccio{7}{Triángulos}

\ex{1} Clasifica según los lados los triángulos que tienen estas medidas (en cm):\begin{tasks}(4)
  \task 5, 5, 5
  \task 4, 7, 7
  \task 3, 5, 6
  \task 8, 8, 3
\end{tasks}\recorda{Según los lados: \textbf{equilátero} (3 iguales), \textbf{isósceles} (2 iguales), \textbf{escaleno} (ninguno igual). Según los ángulos: \textbf{acutángulo}, \textbf{rectángulo} u \textbf{obtusángulo}. Los tres ángulos suman $180^\circ$.}

\ex{1} Dos ángulos de un triángulo miden lo que se indica. Calcula el tercero y clasifica el triángulo según los ángulos:\begin{tasks}(4)
  \task $40^\circ$ y $60^\circ$
  \task $90^\circ$ y $35^\circ$
  \task $120^\circ$ y $25^\circ$
  \task $45^\circ$ y $45^\circ$
\end{tasks}

\ex{1} ¿Se puede construir un triángulo con segmentos de estas longitudes (en cm)?\begin{tasks}(4)
  \task 3, 4, 5
  \task 2, 3, 6
  \task 5, 5, 9
  \task 4, 4, 8
\end{tasks}

\ex{2} Calcula:\begin{tasks}(1)
  \task Los ángulos de la base de un triángulo isósceles en el que el ángulo desigual mide $50^\circ$.
  \task El ángulo desigual de un triángulo isósceles en el que cada ángulo de la base mide $72^\circ$.
  \task Los ángulos agudos de un triángulo rectángulo isósceles.
\end{tasks}

\ex{2} ¿Puede existir un triángulo que sea…\begin{tasks}(2)
  \task rectángulo e isósceles?
  \task obtusángulo y equilátero?
  \task acutángulo y escaleno?
  \task rectángulo y equilátero?
\end{tasks}

\ex{2} Un ángulo agudo de un triángulo rectángulo mide $27^\circ\,35'$. ¿Cuánto mide el otro ángulo agudo?

\ex{2} ¿Qué punto notable de un triángulo es…\begin{tasks}(1)
  \task el centro de la circunferencia que pasa por los tres vértices?
  \task el punto donde se cortan las tres alturas?
  \task el centro de gravedad del triángulo?
  \task el punto que está a la misma distancia de los tres lados?
\end{tasks}

\ex{3} Los ángulos de un triángulo miden $x$, $x+15^\circ$ y $x+30^\circ$. Calcúlalos y clasifica el triángulo.

\ex{3} Dos lados de un triángulo miden 5\,cm y 12\,cm. ¿Qué longitudes (en centímetros enteros) puede tener el tercer lado?

\ex{3} Prolongamos un lado de un triángulo y, en uno de los vértices, el ángulo adyacente al interior que se forma fuera del triángulo mide $130^\circ$. Otro ángulo del triángulo mide $60^\circ$. Calcula los tres ángulos del triángulo.


% =====================================================
\seccio{8}{Cuadriláteros}

\ex{1} Escribe el nombre de cada cuadrilátero (las marcas indican lados iguales):\begin{tasks}(3)
  \task \begin{tikzpicture}[x=0.5cm,y=0.5cm,baseline=(current bounding box.center),line join=round,line cap=round]
\draw[fill=teal!12,draw=teal,thick] (0.000,0.000) -- (2.000,0.000) -- (2.000,2.000) -- (0.000,2.000) -- cycle;
\draw[thin] (0.000,0.440) -- (0.440,0.440) -- (0.440,0.000);
\draw[thin] (1.560,0.000) -- (1.560,0.440) -- (2.000,0.440);
\draw[thin] (2.000,1.560) -- (1.560,1.560) -- (1.560,2.000);
\draw[thin] (0.440,2.000) -- (0.440,1.560) -- (0.000,1.560);
\draw[thick] (1.000,0.240) -- (1.000,-0.240);
\draw[thick] (1.760,1.000) -- (2.240,1.000);
\draw[thick] (1.000,1.760) -- (1.000,2.240);
\draw[thick] (0.240,1.000) -- (-0.240,1.000);
\end{tikzpicture}
  \task \begin{tikzpicture}[x=0.5cm,y=0.5cm,baseline=(current bounding box.center),line join=round,line cap=round]
\draw[fill=teal!12,draw=teal,thick] (0.000,0.000) -- (1.600,-1.100) -- (3.200,0.000) -- (1.600,1.100) -- cycle;
\draw[thick] (0.936,-0.352) -- (0.664,-0.748);
\draw[thick] (2.264,-0.352) -- (2.536,-0.748);
\draw[thick] (2.264,0.352) -- (2.536,0.748);
\draw[thick] (0.936,0.352) -- (0.664,0.748);
\end{tikzpicture}
  \task \begin{tikzpicture}[x=0.5cm,y=0.5cm,baseline=(current bounding box.center),line join=round,line cap=round]
\draw[fill=teal!12,draw=teal,thick] (0.000,0.000) -- (3.400,0.000) -- (2.500,1.700) -- (0.900,1.700) -- cycle;
\draw[thick] (2.738,0.738) -- (3.162,0.962);
\draw[thick] (0.662,0.738) -- (0.238,0.962);
\end{tikzpicture}
  \task \begin{tikzpicture}[x=0.5cm,y=0.5cm,baseline=(current bounding box.center),line join=round,line cap=round]
\draw[fill=teal!12,draw=teal,thick] (0.000,0.000) -- (3.000,0.000) -- (3.000,1.700) -- (0.000,1.700) -- cycle;
\draw[thin] (0.000,0.440) -- (0.440,0.440) -- (0.440,0.000);
\draw[thin] (2.560,0.000) -- (2.560,0.440) -- (3.000,0.440);
\draw[thin] (3.000,1.260) -- (2.560,1.260) -- (2.560,1.700);
\draw[thin] (0.440,1.700) -- (0.440,1.260) -- (0.000,1.260);
\draw[thick] (1.500,0.240) -- (1.500,-0.240);
\draw[thick] (2.760,0.760) -- (3.240,0.760);
\draw[thick] (2.760,0.940) -- (3.240,0.940);
\draw[thick] (1.500,1.460) -- (1.500,1.940);
\draw[thick] (0.240,0.940) -- (-0.240,0.940);
\draw[thick] (0.240,0.760) -- (-0.240,0.760);
\end{tikzpicture}
  \task \begin{tikzpicture}[x=0.5cm,y=0.5cm,baseline=(current bounding box.center),line join=round,line cap=round]
\draw[fill=teal!12,draw=teal,thick] (0.000,0.000) -- (2.600,0.000) -- (3.500,1.600) -- (0.900,1.600) -- cycle;
\draw[thick] (1.300,0.240) -- (1.300,-0.240);
\draw[thick] (2.797,0.839) -- (3.215,0.604);
\draw[thick] (2.885,0.996) -- (3.303,0.761);
\draw[thick] (2.200,1.360) -- (2.200,1.840);
\draw[thick] (0.703,0.761) -- (0.285,0.996);
\draw[thick] (0.615,0.604) -- (0.197,0.839);
\end{tikzpicture}
  \task \begin{tikzpicture}[x=0.5cm,y=0.5cm,baseline=(current bounding box.center),line join=round,line cap=round]
\draw[fill=teal!12,draw=teal,thick] (0.000,0.000) -- (3.200,0.000) -- (2.000,1.800) -- (0.000,1.800) -- cycle;
\draw[thin] (0.000,0.440) -- (0.440,0.440) -- (0.440,0.000);
\draw[thin] (0.440,1.800) -- (0.440,1.360) -- (0.000,1.360);
\end{tikzpicture}
\end{tasks}\recorda{\textbf{Paralelogramos} (lados opuestos paralelos): cuadrado, rectángulo, rombo y romboide. \textbf{Trapecios} (solo dos lados paralelos): rectángulo, isósceles y escaleno. \textbf{Trapezoides}: ningún lado paralelo.}

\ex{1} Tres ángulos de un cuadrilátero miden lo que se indica. Calcula el cuarto:\begin{tasks}(3)
  \task $90^\circ$, $90^\circ$, $100^\circ$
  \task $75^\circ$, $105^\circ$, $75^\circ$
  \task $120^\circ$, $60^\circ$, $95^\circ$
\end{tasks}\recorda{Los cuatro ángulos de un cuadrilátero suman $360^\circ$.}

\ex{1} ¿Qué paralelogramos tienen…\begin{tasks}(2)
  \task las diagonales iguales?
  \task las diagonales perpendiculares?
  \task los cuatro ángulos rectos?
  \task los cuatro lados iguales?
\end{tasks}

\ex{2} Calcula los ángulos que faltan:\begin{tasks}(1)
  \task Un rombo con un ángulo de $54^\circ$.
  \task Un romboide con un ángulo de $112^\circ$.
  \task Un trapecio rectángulo con un ángulo agudo de $63^\circ$.
\end{tasks}

\ex{2} Di si es cierto o falso:\begin{tasks}(1)
  \task Todo cuadrado es un rectángulo.
  \task Todo rombo es un cuadrado.
  \task Todo rectángulo es un paralelogramo.
  \task Un trapecio puede tener tres ángulos rectos.
\end{tasks}

\ex{3} Responde:\begin{tasks}(1)
  \task Los ángulos de un cuadrilátero miden $x$, $2x$, $3x$ y $4x$. ¿Cuánto mide cada uno?
  \task En un romboide, un ángulo es el triple de su consecutivo. ¿Cuánto miden los cuatro ángulos?
\end{tasks}


% =====================================================
\seccio{9}{Ángulos de los polígonos. Polígonos regulares}

\ex{1} ¿Cuánto suman los ángulos interiores de cada polígono?\begin{tasks}(5)
  \task pentágono
  \task hexágono
  \task octógono
  \task decágono
  \task dodecágono
\end{tasks}\recorda{Los ángulos de un polígono de $n$ lados suman $(n-2)\cdot180^\circ$. En un polígono regular, cada ángulo interior mide $\dfrac{(n-2)\cdot180^\circ}{n}$ y el ángulo central mide $\dfrac{360^\circ}{n}$.}

\ex{1} Calcula el ángulo interior y el ángulo central de estos polígonos regulares:\begin{tasks}(3)
  \task triángulo equilátero
  \task cuadrado
  \task pentágono
  \task hexágono
  \task octógono
  \task decágono
\end{tasks}

\ex{2} El ángulo central de un polígono regular mide $24^\circ$. ¿Cuántos lados tiene? ¿Cuánto mide cada ángulo interior?

\ex{2} Cada ángulo interior de un polígono regular mide $160^\circ$. ¿Cuántos lados tiene?

\ex{2} Los ángulos interiores de un polígono suman $2340^\circ$. ¿Cuántos lados tiene?

\ex{3} ¿Puede haber un polígono regular con ángulos interiores de $100^\circ$? ¿Y de $150^\circ$?

\ex{3} ¿Con cuáles de estos polígonos regulares se puede cubrir el suelo sin dejar huecos: triángulos, cuadrados, pentágonos, hexágonos u octógonos? ¿Y combinando octógonos y cuadrados?

\ex{3} Unimos el centro de un hexágono regular con los seis vértices. ¿Qué tipo de triángulos se obtienen? Razónalo con los ángulos.

\vfill
\begin{center}\color{muted}\small\faCheckCircle[regular]\ \ ¡Ánimo! Si te cuesta algún ejercicio, repasa los apuntes interactivos de la unidad.\end{center}

\end{document}
